% =====================================================================
\chapter{Objetivos}
\label{cap:objetivos}

% ---------------------------------------------------------------------
\section{Problemática observada}

En la conversación cotidiana sobre la bolsa argentina está muy extendida la idea de que «las
noticias mueven las acciones»: un titular negativo sobre el dólar o sobre una empresa
precedería a una caída de su cotización, y una buena noticia, a una suba. La observación tiene
sustento intuitivo, porque los días de noticias grandes, como los posteriores a una elección,
son también días de movimientos fuertes del S\&P~Merval. Pero esa coincidencia no dice qué
vino primero. Es posible que las noticias muevan al mercado, y también que los medios cuenten
lo que el mercado ya hizo, con un tono que acompaña la suba o la baja. El problema de este
trabajo es separar esas dos posibilidades con datos: ordenar en el tiempo el tono y la
cantidad de las noticias económicas y el precio y el volumen de las acciones, y ver cuál
anticipa a cuál.

% ---------------------------------------------------------------------
\section{Hipótesis de investigación}

El plan de trabajo define el \textit{shock informativo} como un aumento anormal del volumen
de noticias publicadas por medios digitales argentinos sobre temas económico-financieros, y
plantea las siguientes hipótesis:

\begin{quote}
\textbf{Hipótesis alternativa (H1).} «El shock informacional tiene un impacto medible y
estadísticamente significativo sobre el comportamiento bursátil —precio y/o volumen de
transacciones— de las empresas argentinas que cotizan en bolsa.»

\textbf{Hipótesis nula (H0).} «El shock informacional no tiene un impacto estadísticamente
significativo sobre el comportamiento bursátil de las empresas argentinas.»
\end{quote}

Para contrastarlas, el shock se midió de dos maneras complementarias: por su \textit{tono}
(días en que las noticias son mucho más positivas o negativas que lo habitual) y por su
\textit{cantidad} (días con muchas más notas que lo habitual, la definición literal del plan).
La dirección que interesa es noticias \flecha{} mercado: que el shock preceda al movimiento
del precio o del volumen. Si los datos muestran la dirección contraria, o ninguna, la
hipótesis no se sostiene, y ese también es un resultado.

% ---------------------------------------------------------------------
\section{Objetivo general}

\begin{itemize}
  \item Modelar y cuantificar el impacto de shocks informativos provenientes de portales de
        noticias argentinos sobre el comportamiento bursátil —precio y volumen de
        transacciones— de empresas argentinas que cotizan en bolsa.
\end{itemize}

% ---------------------------------------------------------------------
\section{Objetivos específicos}

\begin{enumerate}
  \item Construir y consolidar un corpus de noticias de economía y finanzas de los
        principales portales digitales argentinos, estructurado como serie temporal por
        portal y fecha.
  \item Diseñar e implementar un modelo de análisis de sentimiento financiero en español,
        adaptando la arquitectura FinBERT mediante el ajuste fino de BETO sobre noticias
        financieras argentinas.
  \item Integrar datos de \textit{trending topics} de X/Twitter como variable complementaria
        para la detección y validación de shocks informativos.
  \item Definir y cuantificar variables de shock informativo: volumetría de publicaciones,
        temperatura del sentimiento agregado y presencia de entidades en el flujo noticioso.
  \item Agrupar las empresas del mercado accionario argentino en conjuntos significativos,
        considerando la clasificación sectorial y la sensibilidad al flujo informativo.
  \item Modelar la relación entre las variables de shock informativo y las variables
        bursátiles, incorporando desfases temporales para capturar el impacto diferido del
        flujo de noticias.
\end{enumerate}

La tabla~\ref{tab:objetivos} resume cómo se cumplió cada objetivo y en qué parte del
documento se desarrolla. El objetivo~3 no se realizó: los datos de X/Twitter no llegaron a
integrarse al análisis y quedan como trabajo futuro. En el objetivo~6, el plan preveía también
un modelo de rezagos distribuidos (DLM); se cumplió con un VAR, que ya incorpora los efectos
rezagados, y el DLM queda como línea futura.

\begin{table}[H]
\centering\small
\begin{tabularx}{\linewidth}{@{}C{0.6cm}YL{1.9cm}L{2.3cm}@{}}
\toprule
\textbf{Obj.} & \textbf{Qué se hizo} & \textbf{Estado} & \textbf{Dónde} \\
\midrule
1 & 14.012 notas de 7 portales, depuradas y asignadas a ruedas con la regla de las 17:00 & Cumplido & §\ref{sec:recoleccion}, §\ref{sec:alineacion} \\
2 & FinBETO, 4 clases; F1 macro de 0,685 en noticias no vistas, por encima del modelo de referencia & Cumplido & §\ref{sec:finbeto-modelo} \\
3 & Los datos de X/Twitter no se integraron & No realizado & §\ref{sec:futuro} \\
4 & Atención anormal (cantidad de notas), tono anormal, ruedas de tono extremo, tono propio de cada empresa y shock de cantidad según el plan & Cumplido & §\ref{sec:tono-rueda}, §\ref{sec:res-u7b} \\
5 & Clustering de las 47 acciones y análisis por sector y por empresa & Cumplido & §\ref{sec:multivariado} \\
6 & VAR, Granger, impulso-respuesta, estudios de eventos y análisis por régimen; el DLM queda como trabajo futuro & Cumplido (con VAR) & cap.~\ref{cap:resultados} \\
\bottomrule
\end{tabularx}
\caption{Estado de los objetivos específicos. Fuente: elaboración propia.}
\label{tab:objetivos}
\end{table}

% ---------------------------------------------------------------------
\section{Antecedentes}
\label{sec:antecedentes}

\subsection{Noticias y precios: ¿anticipan o reflejan?}

El antecedente central es \citet{tetlock2007}, que midió el pesimismo de una columna diaria
del \textit{Wall Street Journal} y encontró dos cosas: el pesimismo anticipa una baja de los
precios que luego se revierte, y los retornos bajos anticipan más pesimismo en la columna. Es
decir, la relación existe en las dos direcciones. \citet{garcia2013}, con columnas
financieras del \textit{New York Times} entre 1905 y 2005, mostró que el tono anticipa los
retornos casi solo en las recesiones, lo que sugiere que el efecto depende del estado del
mercado. Otros trabajos separaron el efecto de la cobertura del de su contenido:
\citet{fang2009} encontraron que las acciones sin cobertura de medios rinden más que las muy
cubiertas, \citet{engelberg2011} usaron la cobertura local de un mismo hecho para identificar
un efecto causal de los medios sobre el volumen operado, y \citet{dougal2012} mostraron que la
identidad del periodista que escribe la columna predice el retorno del día siguiente.
\citet{da2011} propusieron medir la atención de los inversores con las búsquedas en internet,
una variable de cantidad y no de tono.

Trabajos más recientes retoman la pregunta de fondo de esta tesina. \citet{verma2026} usan el
shock de 2020 para preguntar si las noticias influyen en el precio o reflejan información que
el precio ya tiene; \citet{jazbec2020} miden cuánta información pasa de las noticias de la web
a los precios con medidas de transferencia de información; \citet{kargarzadeh2026} estudian en
qué momento el precio incorpora una noticia —antes, el mismo día o después—, y
\citet{glasserman2025} se preguntan si las noticias publicadas de noche explican el retorno
entre el cierre y la apertura. \citet{knuteson2020} documenta que el retorno nocturno y el del
día se comportan de forma muy distinta, lo que justifica separarlos. Sobre la cantidad de
noticias, \citet{cao2025hype} construyen un índice de atención mediática por acción, y
\citet{guo2017} muestran que dos empresas que aparecen en una misma noticia comparten
atención y movimientos de precio. \citet{ardia2022} usan el \textit{tono anormal} de los
medios —el tono comparado con lo habitual— en estudios de eventos alrededor de anuncios de
resultados, y \citet{chen2025topics} critican resumir el sentimiento de un día en un solo
número sin separar los temas.

\subsection{Medición del sentimiento con procesamiento de lenguaje natural}

La medición del tono pasó de los diccionarios de palabras a los modelos de lenguaje. El
diccionario financiero de \citet{loughran2011} mostró que las listas de palabras generales
clasifican mal el lenguaje financiero (por ejemplo, «pasivo» no es negativo en un balance).
\citet{malo2014phrasebank} publicaron el \textit{Financial PhraseBank}, un conjunto de
oraciones financieras etiquetadas por personas, y \citet{araci2019finbert} ajustó sobre él un
modelo BERT \citep{devlin2019bert} para clasificar sentimiento financiero en inglés
(FinBERT). Para el español, \citet{canete2020beto} entrenaron BETO, el BERT en español de la
Universidad de Chile, y \citet{gutierrez2022maria} la familia MarIA del Barcelona
Supercomputing Center. \citet{perez2021pysentimiento} publicaron pysentimiento, con modelos de
sentimiento en español entrenados con tuits. En el dominio financiero en español, la
competencia FinancES 2023 \citep{marin2023finances} y los trabajos de \citet{pan2023peerj},
\citet{garciadiaz2023ieee} y \citet{rodriguezgarcia2023} evaluaron modelos tipo BERT sobre
titulares y tuits económicos; en la comparación de \citet{pan2023peerj}, BETO y MarIA son los
de mejor desempeño.

Sobre cómo pasar de oraciones a series de tiempo, \citet{mavillonio2026} construyen índices
diarios de sentimiento con modelos tipo BERT frase por frase y los comparan con los de
diccionario; \citet{stan2026} y \citet{carbo2026} muestran que la serie es ruidosa cuando hay
días con pocas noticias y proponen darles menos peso, y \citet{correarocha2026} documenta que
el tono de las noticias económicas es muy persistente. A nivel de empresa,
\citet{tang2023finentity} y \citet{sinha2023sentfin} clasifican el sentimiento hacia cada
entidad nombrada en una noticia, ya que una misma nota puede ser buena para una empresa y mala
para otra.

\subsection{Argentina y Latinoamérica}

En la región los antecedentes son pocos. \citet{brana2016} clasificaron frases de noticias
financieras argentinas como positivas, negativas o neutras, sin cruzarlas con precios.
\citet{galdi2018} midieron el tono de 1.237 ediciones del diario \textit{Valor Econômico} y
encontraron que las palabras de incertidumbre se asocian con menor rentabilidad de la bolsa
brasileña. \citet{cruzmontana2025} cruzaron el tono de noticias financieras con el índice
COLCAP de Colombia. \citet{coronado2022,coronado2023} probaron la causalidad de Granger
desde índices de noticias sobre el COVID-19 hacia las bolsas latinoamericanas, Argentina
incluida, pero con cantidad de noticias y no con tono. \citet{rojas2015} estudiaron la
causalidad entre retorno y volumen en las bolsas de la región, y \citet{chniguir2022} la
causalidad entre sentimiento y retornos en Marruecos y Túnez, dos mercados emergentes chicos
como el argentino. \citet{mukherjee2025} y \citet{wu2026} aplican pruebas de Granger y redes
de contagio del sentimiento entre noticias y redes sociales. Para Argentina,
\citet{calandroni2026} estudió el efecto de las elecciones sobre el Merval entre 2011 y 2023
con un estudio de eventos, y \citet{amburgey2025} mostró, para Estados Unidos, que un mismo
shock electoral mueve a los sectores en direcciones distintas.

\subsection{El vacío que cubre este trabajo}

No se encontró un trabajo que cruce el tono de la prensa económica argentina con el precio y
el volumen de las acciones de BYMA. Los antecedentes en español clasifican el sentimiento sin
cruzarlo con el mercado; los latinoamericanos que sí lo cruzan son de Brasil y Colombia, o
usan cantidad de noticias y no tono. Esta tesina aporta ese cruce para Argentina, con un
modelo de sentimiento propio, tono por empresa, pruebas en las dos direcciones y una prueba
que usa la hora de publicación de cada nota para ordenar noticias y precios en el tiempo.

% ---------------------------------------------------------------------
\section{Metodología}
\label{sec:metodologia}

\textbf{Tipo de estudio.} Es un estudio observacional, longitudinal y correlacional: no
interviene sobre las variables, las sigue en el tiempo (244 ruedas) y analiza su relación y su
orden temporal. La población es el conjunto de acciones de empresas argentinas que cotizan en
BYMA; la muestra son 47 acciones que cumplen criterios de cobertura y liquidez, 17 de ellas
con noticias suficientes para medir su tono propio. El corpus son las notas de las secciones
de economía y finanzas de siete portales nacionales.

El trabajo se organizó como un \textit{pipeline} de tres partes, siguiendo las fases del
proceso CRISP-DM \citep{wirth2000crispdm}:

\begin{itemize}
  \item \textbf{Lado bursátil.} Selección del panel de acciones a partir de Yahoo Finance,
        clustering, volatilidad realizada del mercado (VIX argentino), pruebas de
        estacionariedad, volatilidad por acción y régimen de mercado con un modelo de cambio
        de régimen de Markov.
  \item \textbf{Lado de noticias.} Recolección del corpus (Wayback Machine y mapas del sitio),
        depuración, etiquetado de oraciones con un modelo de lenguaje, entrenamiento de
        FinBETO, serie diaria de sentimiento y reconocimiento de entidades (NER) para medir el
        tono de cada empresa.
  \item \textbf{Unión.} Alineación de cada nota con la rueda en la que el mercado pudo
        reaccionar (regla de las 17:00), series de tono y atención por rueda y por empresa,
        datasets unificados, estacionariedad, VAR con pruebas de Granger e impulso-respuesta,
        prueba con las notas publicadas con el mercado cerrado, estudios de eventos y
        robustez.
\end{itemize}

Dos especificaciones (la prueba pre-apertura y el estudio de eventos del shock de cantidad)
se escribieron y guardaron antes de correr los cálculos, y todas las familias de pruebas se
corrigieron por comparaciones múltiples. El detalle de cada paso está en los
capítulos~\ref{cap:tecnicas} y~\ref{cap:desarrollo}.

% ---------------------------------------------------------------------
\section{Impacto esperado}

\textbf{Científico.} El trabajo aporta evidencia sobre la relación entre medios y mercado en un
mercado emergente y en español, donde casi no la hay, y un conjunto de oraciones de prensa
económica argentina etiquetadas por sentimiento, que no existía. El método es replicable:
cada paso tiene su script, sus controles automáticos y su documentación.

\textbf{Práctico.} Para inversores y analistas, el resultado indica si el tono de los medios
es una señal que se adelanta al mercado o un reflejo de él. Para reguladores, ofrece una forma
de medir cómo la cobertura mediática acompaña la actividad bursátil.

\textbf{Metodológico.} El \textit{pipeline} es reproducible de punta a punta: un solo comando
rehace toda la unión y verifica que los 64 archivos de resultados salgan idénticos y que los
66 controles automáticos pasen.

\textbf{Por qué de a dos.} La tesina combina dos campos con técnicas propias: el
procesamiento de lenguaje natural (recolectar y depurar miles de notas, entrenar y evaluar un
modelo de sentimiento, reconocer entidades) y la econometría financiera (armar y depurar un
panel de acciones, medir volatilidad y regímenes, modelar series de tiempo). La carga de
trabajo de cada lado justifica un equipo de dos. La división fue la siguiente: el lado de
noticias (corpus, sentimiento y NER) estuvo a cargo de Gino Frigoni; el lado bursátil
(panel, clustering, volatilidad, estacionariedad y régimen), de Nicolás Diddi. La unión de
los dos lados, que responde la pregunta de la tesina, integra el trabajo de ambos.

% ---------------------------------------------------------------------
\section{Cronograma}

El plan de trabajo se presentó en mayo de 2026. A partir de entonces el trabajo avanzó por
lados: cada uno cerró sus resultados con su propia documentación y recién después se unieron.
La tabla~\ref{tab:cronograma} registra los hitos documentados del proceso.

\begin{table}[H]
\centering\small
\begin{tabularx}{\linewidth}{@{}L{3.4cm}YL{2.6cm}@{}}
\toprule
\textbf{Etapa} & \textbf{Hito} & \textbf{Fecha} \\
\midrule
Plan de trabajo & Presentación del proyecto (hipótesis, objetivos y metodología) & Mayo de 2026 \\
Corpus de noticias & Recolección de notas publicadas hasta el 3 de mayo de 2026 & Hasta mayo de 2026 \\
Lado bursátil & Panel de 47 acciones, VIX, ADF, volatilidad y régimen; revisión cerrada & Octubre de 2026 \\
Lado de noticias & Auditoría del corpus, nueva medición de FinBETO y cruce con el NER; versión final & 18/10/2026 \\
Unión & Alineación, datasets, VAR y Granger, estudios de eventos y robustez (U1 a U9) & 17 al 24/10/2026 \\
Actualización del plan & Estado de los objetivos y corrección del apellido de un autor (versión 3) & Octubre de 2026 \\
Informe final & Redacción de esta tesina & Octubre de 2026 \\
\bottomrule
\end{tabularx}
\caption{Hitos del trabajo. Fuente: elaboración propia a partir de la documentación de cada
etapa.}
\label{tab:cronograma}
\end{table}

Respecto de lo planificado hubo tres cambios. Los datos de X/Twitter no se integraron
(objetivo~3). El DLM se reemplazó por el VAR, que ya trabaja con rezagos. Y se sumaron
pruebas que el plan no preveía: el tono propio de cada empresa, la prueba con las notas
publicadas con el mercado cerrado y las doce variantes de robustez.
